% =====================================================================
% Preâmbulo do mapeamento sistemático
% Inclua no preâmbulo do documento principal:  % =====================================================================
% Preâmbulo do mapeamento sistemático
% Inclua no preâmbulo do documento principal:  % =====================================================================
% Preâmbulo do mapeamento sistemático
% Inclua no preâmbulo do documento principal:  % =====================================================================
% Preâmbulo do mapeamento sistemático
% Inclua no preâmbulo do documento principal:  \input{preambulo_mapeamento}
% (se a pasta estiver em outro lugar, ajuste o caminho e o \mappath abaixo)
% =====================================================================
\usepackage{graphicx}
\usepackage{booktabs}
\usepackage{longtable}
\usepackage{multirow}
\usepackage{array}

% Caminho da pasta do mapeamento em relação ao arquivo principal.
% Ex.: se a pasta se chama "mapeamento_latex" e está ao lado do main.tex:
%      \renewcommand{\mappath}{mapeamento_latex/}
\providecommand{\mappath}{}

% Citação textual: abnTeX2 usa \citeonline; com outros pacotes cai em \cite.
\providecommand{\citeonline}[1]{\cite{#1}}

% Linha de fonte abaixo de figuras/tabelas/quadros (padrão ABNT).
\providecommand{\mapfonte}[1]{\par\smallskip{\footnotesize #1\par}}

% ---------------------------------------------------------------------
% Quadros
% Se a sua classe/template já define o ambiente/contador "quadro"
% (comum em templates abnTeX2), ele é reaproveitado. Caso contrário,
% um contador é criado. \mapquadro{legenda}{rótulo} escreve a legenda
% "Quadro N -- ..." e registra o item na lista de quadros (arquivo .loq).
% ---------------------------------------------------------------------
\providecommand{\quadroname}{Quadro}
\providecommand{\quadroautorefname}{Quadro}
% \autoref vem do hyperref (carregado pelo abnTeX2). Sem hyperref, vira \ref.
\AtBeginDocument{\providecommand{\autoref}[1]{\ref{#1}}}
\makeatletter
\@ifundefined{c@quadro}{\newcounter{quadro}}{}
\newcommand{\mapquadro}[2]{%
  \par\addvspace{\medskipamount}%
  \refstepcounter{quadro}\label{#2}%
  \addcontentsline{loq}{quadro}{\protect\numberline{\thequadro}#1}%
  \nopagebreak{\centering\small \quadroname~\thequadro~--~#1\par}\nopagebreak[4]\vspace{2pt}\nopagebreak[4]}
\makeatother

% (se a pasta estiver em outro lugar, ajuste o caminho e o \mappath abaixo)
% =====================================================================
\usepackage{graphicx}
\usepackage{booktabs}
\usepackage{longtable}
\usepackage{multirow}
\usepackage{array}

% Caminho da pasta do mapeamento em relação ao arquivo principal.
% Ex.: se a pasta se chama "mapeamento_latex" e está ao lado do main.tex:
%      \renewcommand{\mappath}{mapeamento_latex/}
\providecommand{\mappath}{}

% Citação textual: abnTeX2 usa \citeonline; com outros pacotes cai em \cite.
\providecommand{\citeonline}[1]{\cite{#1}}

% Linha de fonte abaixo de figuras/tabelas/quadros (padrão ABNT).
\providecommand{\mapfonte}[1]{\par\smallskip{\footnotesize #1\par}}

% ---------------------------------------------------------------------
% Quadros
% Se a sua classe/template já define o ambiente/contador "quadro"
% (comum em templates abnTeX2), ele é reaproveitado. Caso contrário,
% um contador é criado. \mapquadro{legenda}{rótulo} escreve a legenda
% "Quadro N -- ..." e registra o item na lista de quadros (arquivo .loq).
% ---------------------------------------------------------------------
\providecommand{\quadroname}{Quadro}
\providecommand{\quadroautorefname}{Quadro}
% \autoref vem do hyperref (carregado pelo abnTeX2). Sem hyperref, vira \ref.
\AtBeginDocument{\providecommand{\autoref}[1]{\ref{#1}}}
\makeatletter
\@ifundefined{c@quadro}{\newcounter{quadro}}{}
\newcommand{\mapquadro}[2]{%
  \par\addvspace{\medskipamount}%
  \refstepcounter{quadro}\label{#2}%
  \addcontentsline{loq}{quadro}{\protect\numberline{\thequadro}#1}%
  \nopagebreak{\centering\small \quadroname~\thequadro~--~#1\par}\nopagebreak[4]\vspace{2pt}\nopagebreak[4]}
\makeatother

% (se a pasta estiver em outro lugar, ajuste o caminho e o \mappath abaixo)
% =====================================================================
\usepackage{graphicx}
\usepackage{booktabs}
\usepackage{longtable}
\usepackage{multirow}
\usepackage{array}

% Caminho da pasta do mapeamento em relação ao arquivo principal.
% Ex.: se a pasta se chama "mapeamento_latex" e está ao lado do main.tex:
%      \renewcommand{\mappath}{mapeamento_latex/}
\providecommand{\mappath}{}

% Citação textual: abnTeX2 usa \citeonline; com outros pacotes cai em \cite.
\providecommand{\citeonline}[1]{\cite{#1}}

% Linha de fonte abaixo de figuras/tabelas/quadros (padrão ABNT).
\providecommand{\mapfonte}[1]{\par\smallskip{\footnotesize #1\par}}

% ---------------------------------------------------------------------
% Quadros
% Se a sua classe/template já define o ambiente/contador "quadro"
% (comum em templates abnTeX2), ele é reaproveitado. Caso contrário,
% um contador é criado. \mapquadro{legenda}{rótulo} escreve a legenda
% "Quadro N -- ..." e registra o item na lista de quadros (arquivo .loq).
% ---------------------------------------------------------------------
\providecommand{\quadroname}{Quadro}
\providecommand{\quadroautorefname}{Quadro}
% \autoref vem do hyperref (carregado pelo abnTeX2). Sem hyperref, vira \ref.
\AtBeginDocument{\providecommand{\autoref}[1]{\ref{#1}}}
\makeatletter
\@ifundefined{c@quadro}{\newcounter{quadro}}{}
\newcommand{\mapquadro}[2]{%
  \par\addvspace{\medskipamount}%
  \refstepcounter{quadro}\label{#2}%
  \addcontentsline{loq}{quadro}{\protect\numberline{\thequadro}#1}%
  \nopagebreak{\centering\small \quadroname~\thequadro~--~#1\par}\nopagebreak[4]\vspace{2pt}\nopagebreak[4]}
\makeatother

% (se a pasta estiver em outro lugar, ajuste o caminho e o \mappath abaixo)
% =====================================================================
\usepackage{graphicx}
\usepackage{booktabs}
\usepackage{longtable}
\usepackage{multirow}
\usepackage{array}

% Caminho da pasta do mapeamento em relação ao arquivo principal.
% Ex.: se a pasta se chama "mapeamento_latex" e está ao lado do main.tex:
%      \renewcommand{\mappath}{mapeamento_latex/}
\providecommand{\mappath}{}

% Citação textual: abnTeX2 usa \citeonline; com outros pacotes cai em \cite.
\providecommand{\citeonline}[1]{\cite{#1}}

% Linha de fonte abaixo de figuras/tabelas/quadros (padrão ABNT).
\providecommand{\mapfonte}[1]{\par\smallskip{\footnotesize #1\par}}

% ---------------------------------------------------------------------
% Quadros
% Se a sua classe/template já define o ambiente/contador "quadro"
% (comum em templates abnTeX2), ele é reaproveitado. Caso contrário,
% um contador é criado. \mapquadro{legenda}{rótulo} escreve a legenda
% "Quadro N -- ..." e registra o item na lista de quadros (arquivo .loq).
% ---------------------------------------------------------------------
\providecommand{\quadroname}{Quadro}
\providecommand{\quadroautorefname}{Quadro}
% \autoref vem do hyperref (carregado pelo abnTeX2). Sem hyperref, vira \ref.
\AtBeginDocument{\providecommand{\autoref}[1]{\ref{#1}}}
\makeatletter
\@ifundefined{c@quadro}{\newcounter{quadro}}{}
\newcommand{\mapquadro}[2]{%
  \par\addvspace{\medskipamount}%
  \refstepcounter{quadro}\label{#2}%
  \addcontentsline{loq}{quadro}{\protect\numberline{\thequadro}#1}%
  \nopagebreak{\centering\small \quadroname~\thequadro~--~#1\par}\nopagebreak[4]\vspace{2pt}\nopagebreak[4]}
\makeatother
